\subsection{Camera-Based Hand Tracking with MediaPipe}
\label{subsec:mediapipe}

\subsubsection{Pipeline Architecture}
MediaPipe Hands \cite{zhang2020mediapipehands} is a two-stage machine learning
pipeline designed for real-time, on-device hand tracking from a single RGB camera.
The first stage is a palm detector that operates on the full input image and locates
palms via an oriented hand bounding box. The second stage is a hand landmark model
that operates on the cropped hand bounding box provided by the palm detector and
returns 21 2.5D landmarks per hand, each described by $x$, $y$, and a relative
depth coordinate defined with respect to the wrist point.

To achieve real-time performance, the pipeline avoids running the palm detector on
every frame. Instead, once the hand has been located, the bounding box for the
current frame is derived directly from the landmark predictions of the previous
frame, so the detector is only reapplied on the first frame or when tracking is
lost. To determine whether tracking has been lost, the hand landmark model outputs
an additional scalar representing the confidence that a reasonably aligned hand is
present in the input crop. When this confidence falls below a predefined threshold,
the palm detector is triggered again to reinitialise tracking.

\subsubsection{Gesture Recognition}
Gesture recognition is built on top of the predicted hand skeleton. The state of
each finger --- for example, whether it is bent or straight --- is determined via
the accumulated angles of its joints. The resulting set of finger states is then
mapped to a set of predefined gestures. Beyond static gesture recognition, a
sequence of landmarks can also be used to predict dynamic gestures.
